% !TEX root = ReportMain.tex
\section*{Introduction}
\addcontentsline{toc}{section}{Introduction}

Le présent rapport rend compte du projet de stage effectué au sein de Yazaki Morocco Kénitra (YMOK), filiale marocaine du groupe Yazaki, équipementier automobile japonais de premier rang. Le projet vise à concevoir et déployer un système de supervision thermique et de maintenance prédictive pour la zone de \textit{foaming} de l'usine, où douze moules sont utilisés pour le surmoulage de cosses électriques.

\paragraph{Contexte industriel.}
Dans le procédé de foaming, chaque moule est maintenu à une température de consigne de \SI{45}{\celsius} par un bain d'eau régulé. Six de ces moules sont instrumentés par des sondes PT100 \SI{100}{\ohm} connectées à un bus RS485 Modbus RTU, et un débitmètre à effet Hall YF-S201 mesure le débit d'eau de refroidissement. L'eau, chauffée par des résistances électriques réparties en quatre groupes, circule en boucle fermée.

\paragraph{Problématique.}
L'eau de chauffage, non adoucie, provoque un dépôt progressif de calcaire (carbonate de calcium) sur les parois internes des moules et sur les sondes de température. Ce dépôt agit comme un isolant thermique : la température mesurée par la sonde diffère de la température réelle du bain, rendant la régulation inefficace et faussant les données de supervision. À terme, l'encrassement peut entraîner des défauts de surmoulage, des arrêts de production pour nettoyage, et une dégradation accélérée des équipements (résistances et pompes).

\paragraph{Solution proposée.}
Pour répondre à cette problématique, le projet déploie une solution en deux volets~:
\begin{itemize}
\item \textbf{Un système de supervision temps réel} fondé sur un Raspberry Pi 4 qui acquiert les mesures de température et de débit, les stocke dans une base de données temporelle InfluxDB, et les diffuse via WebSocket à une interface de visualisation React.
\item \textbf{Un pipeline de diagnostic et de maintenance prédictive} composé de quatre modules~: un modèle physique Grey-Box qui estime l'épaisseur de calcaire par bilan thermique (soft sensing) ; un détecteur d'anomalies non supervisé par Isolation Forest ; un classifieur hybride qui distingue les causes de dérive (règles métier et Random Forest) ; et un régresseur Ridge qui prédit le temps restant avant le prochain seuil d'alerte.
\end{itemize}
En complément, un adoucisseur d'eau est en cours de déploiement pour traiter la cause racine de l'encrassement.

\paragraph{Organisation du rapport.}
Le rapport s'articule en quatre chapitres. Le premier chapitre présente l'entreprise d'accueil, son organisation et ses processus de production. Le deuxième chapitre décrit la problématique technique, l'analyse des modes de défaillance et le cahier des charges. Le troisième chapitre détaille l'architecture et la conception du système, de l'acquisition des données jusqu'aux algorithmes de diagnostic. Le quatrième chapitre expose la réalisation, les tests et les résultats obtenus, et confronte les performances du système au cahier des charges initial.
